\documentclass{uerj}
\usepackage[utf8]{inputenc}
\usepackage[T1]{fontenc}
\usepackage{geometry}
\usepackage{graphicx}
\usepackage{float}
\usepackage{hyperref}
\usepackage[english,indonesian]{babel}

\title{Optimasi Pemrosesan Template LaTeX Menggunakan Pendekatan Berbasis AI}
\author{Academic Support Agent}
\date{\today}

\begin{document}
\maketitle

\begin{abstract}

\end{abstract}

\chapter{PENDAHULUAN}


\chapter{TINJAUAN PUSTAKA}


% Original content (commented out):
% % Modelo de Dissertação em Latex para o PPG em Engenharia Mecânica da UERJ
% Este modelo foi adaptado da versão disponibilizada no site da Engenharia Elétrica da UERJ
% http://www.pel.uerj.br/publico/Modelo_LaTeX_Dissertacao_UERJ.rar
% http://www.pel.uerj.br/defesas/
%
% Utilizei o WinEdt 6.0 com Miktex 2.9
%
% Para gerar o PDF usei a opção PDFTeXfy com o documento [masterthesis.tex] aberto e em foco.
%
% Não consegui usar as figuras em EPS como o modelo original. Usei PNG e JPG sem problemas.
%
% Felipe M. - 20/06/2012
%




\documentclass[a4paper,12pt,oneside,openany]{uerj}

\usepackage[english,brazil]{babel}
%\usepackage[latin1]{inputenc}
\usepackage[utf8]{inputenc}
\usepackage{enumerate}
\usepackage{cite}
\usepackage{epsf,epsfig,psfig}
\usepackage{pagina}
\usepackage{indentfirst}
\usepackage{theorem}
\usepackage{fancyhdr}
\usepackage{setspace}
\usepackage{boxedminipage}
\usepackage{float}
\usepackage{makeidx}
\usepackage{amsmath}
\usepackage[hidelinks]{hyperref}




\makeindex


\newtheorem{deff}{Definição}[section]
\numberwithin{equation}{chapter}

\theoremstyle{plain}

\bibliographystyle{abnt-num}



\begin{document}

\hypersetup{
    colorlinks,
    citecolor=black,
    filecolor=black,
    linkcolor=black,
    urlcolor=black,
    linktoc=all
}

\thispagestyle{empty}\input{./01_Pre_textuais/Capa}
\pagebreak\thispagestyle{empty}\input{./01_Pre_textuais/FolhaDeRosto}
\pagebreak\thispagestyle{empty}\input{./01_Pre_textuais/Fichacatalog}   
\pagebreak\thispagestyle{empty}\input{./01_Pre_textuais/FolhaDeAprovacao}
\pagebreak\thispagestyle{empty}\input{./01_Pre_textuais/Dedicatoria}
\pagebreak\thispagestyle{empty}\input{./01_Pre_textuais/Agradecimento}
%\pagebreak\thispagestyle{empty}\input{./01_Pre_textuais/Epigrafe}    % não coloquei epígrafe no meu trabalho, mas fica aqui a chamada comentada.
\pagebreak\thispagestyle{empty}\input{./01_Pre_textuais/Resumo}
\pagebreak\thispagestyle{empty}\input{./01_Pre_textuais/Abstract}

\fancypagestyle{plain}{
\fancyhf{} % clear all header and footer fields
\renewcommand{\headrulewidth}{0pt}
\renewcommand{\footrulewidth}{0pt}}
\pagestyle{plain}

\pagebreak

\def\listfigurename{LISTA DE FIGURAS}\listoffigures
\def\listtablename{LISTA DE TABELAS}\listoftables
%\input{./01_Pre_textuais/ListaDeSiglas.tex}    % não coloquei LISTA DE SIGLAS no meu trabalho, mas fica aqui a chamada comentada.
\def\contentsname{SUMÁRIO}\tableofcontents

\fancypagestyle{plain}{
\fancyhf{} % clear all header and footer fields
\fancyhead[R]{\thepage}
\setlength{\voffset}{-1cm}
\setlength{\headsep}{1cm}
\renewcommand{\headrulewidth}{0pt}
\renewcommand{\footrulewidth}{0pt}}

\pagestyle{plain}

\pagebreak
\addcontentsline{toc}{chapter}{\hspace{1.7cm}\bfseries INTRODUÇÃO}
\input{./01_Pre_textuais/introducao.tex}



\input{./02_Cap1/exemplo.tex}

% inserir demais capítulos aqui
% -----------------------------
% -----------------------------
% -----------------------------
% -----------------------------





\pagebreak
\addcontentsline{toc}{chapter}{\hspace{1.7cm}\bfseries CONCLUSÃO}
\input{./03_Conclusao/conclusao.tex}

\pagebreak
\addcontentsline{toc}{chapter}{\hspace{1.7cm}\bfseries REFERÊNCIAS}
\def\bibname{REFERÊNCIAS}




% abaixo segue a chamada para o arquivo [.BIB]. Utilizei o programa JABREF para montar o arquivo com minhas referências.
\bibliography{dissertacao}




%felipe% \printindex    %Removi o índice remissivo para a versão oficial do trabalho.


\end{document}

\end{document}